% The Forney-style factor graph of the coupled spatio-sequential model (issue
% #290), drawn from the figure the ticket carries. Input by textbook.tex from
% sec:coupled; kept in its own file because it is a drawing and not prose.
\begin{figure}[htbp]
\centering
\begin{tikzpicture}[
    scale=0.85,
    x={(1.0cm, 0.5cm)},
    y={(0cm, 1.1cm)},
    z={(1.5cm, 0cm)},
    latent/.style={circle, draw, minimum size=0.50cm, inner sep=0.5pt, font=\tiny, fill=white},
    x_latent/.style={circle, draw, minimum size=0.18cm, inner sep=0pt},
    factor_style/.style={rectangle, draw, fill=white, minimum size=0.1cm},
    high_factor_style/.style={regular polygon, regular polygon sides=3, draw, fill=white, minimum size=0.25cm, inner sep=0pt},
    edge_style/.style={draw, thick, black},
    cuboid_edge/.style={draw, thin, black, opacity=0.5},
    dependency/.style={draw, thin, gray, opacity=0.4, bend left=15},
    dependency_inv/.style={draw, thin, gray, opacity=0.4, bend right=15},
    high_dep/.style={draw, thin, black, opacity=1.0},
    colorA/.style={fill=blue!20},
    colorB/.style={fill=red!20}
]
  % Lattice edges and pairwise factors
  \foreach \i in {0,...,4} {
    \foreach \j in {0,...,4} {
      \ifnum\i<4
        \draw[edge_style] (\i, \j, 0) -- (\i+1, \j, 0);
        \node[factor_style] at (\i+0.5, \j, 0) {};
      \fi
      \ifnum\j<4
        \draw[edge_style] (\i, \j, 0) -- (\i, \j+1, 0);
        \node[factor_style] at (\i, \j+0.5, 0) {};
      \fi
    }
  }
  % Markov chain edges and factors, one chain per class
  \foreach \yPos/\xPos/\chainID in {5.5/4/0, -1.5/0/1} {
    \draw[edge_style] (\xPos, \yPos, 2.5) -- (\xPos, \yPos, 3.5);
    \node[factor_style] at (\xPos, \yPos, 3.0) {};
    \foreach \k [count=\t from 0] in {3.5, 4.5, 5.5, 6.5, 7.5} {
       \ifnum\t<4
         \draw[edge_style] (\xPos, \yPos, \k) -- (\xPos, \yPos, \k+1);
         \node[factor_style] at (\xPos, \yPos, \k+0.5) {};
       \fi
    }
  }
  % Observation columns, and the gated factors joining a label to a chain state
  \foreach \i in {0,2,4} {
    \foreach \j in {0,2,4} {
      \draw[dashed, gray!30, thin] (\i, \j, 0) -- (\i, \j, 3.5);
      \foreach \k [count=\t from 0] in {3.5, 4.5, 5.5, 6.5, 7.5} {
        \ifnum\t<4 \draw[cuboid_edge] (\i, \j, \k) -- (\i, \j, \k+1); \fi
        \ifnum\i<4 \draw[cuboid_edge] (\i, \j, \k) -- (\i+2, \j, \k); \fi
        \ifnum\j<4 \draw[cuboid_edge] (\i, \j, \k) -- (\i, \j+2, \k); \fi
        \ifnum\t=0 \ifnum\i=0
            \path[high_dep, bend left=15] (\i, \j, \k) edge node[pos=0.5, high_factor_style] (f-top-\j) {} (4, 5.5, \k);
            \draw[high_dep, bend left=30] (\i, \j, 0) to (f-top-\j);
            \path[high_dep, bend right=15] (\i, \j, \k) edge node[pos=0.5, high_factor_style] (f-bot-\j) {} (0, -1.5, \k);
            \draw[high_dep, bend right=30] (\i, \j, 0) to (f-bot-\j);
        \else
            \path[dependency] (\i, \j, \k) edge (4, 5.5, \k);
            \path[dependency_inv] (\i, \j, \k) edge (0, -1.5, \k);
        \fi \else
            \path[dependency] (\i, \j, \k) edge (4, 5.5, \k);
            \path[dependency_inv] (\i, \j, \k) edge (0, -1.5, \k);
        \fi
      }
    }
  }
  % Label variables on the lattice
  \foreach \i in {0,...,4} {
    \foreach \j in {0,...,4} {
      \ifnum\i=4 \ifnum\j<2 \def\nodeCol{colorA} \else \def\nodeCol{colorB} \fi
      \else \ifnum\i>2 \def\nodeCol{colorB} \else \ifnum\j<\i \def\nodeCol{colorB} \else \def\nodeCol{colorA} \fi \fi \fi
      \node[latent, \nodeCol] at (\i, \j, 0) {$\ell_{\i,\j}$};
    }
  }
  % Chain variables
  \foreach \yPos/\xPos/\chainID in {5.5/4/0, -1.5/0/1} {
    \ifnum\chainID=0 \def\chainCol{colorA} \else \def\chainCol{colorB} \fi
    \node[latent, \chainCol] at (\xPos, \yPos, 2.5) {$\Pi_{\chainID}$};
    \foreach \k [count=\t from 0] in {3.5, 4.5, 5.5, 6.5, 7.5} {
      \node[latent, \chainCol] at (\xPos, \yPos, \k) {$k_{\t,\chainID}$};
    }
  }
  % Observations
  \foreach \i in {0,2,4} {
    \foreach \j in {0,2,4} {
      \foreach \k [count=\t from 0] in {3.5, 4.5, 5.5, 6.5, 7.5} {
        \ifnum\i=4 \ifnum\j<2 \def\xCol{colorA} \else \def\xCol{colorB} \fi
        \else \ifnum\i>2 \def\xCol{colorB} \else \ifnum\j<\i \def\xCol{colorB} \else \def\xCol{colorA} \fi \fi \fi
        \node[x_latent, \xCol] at (\i, \j, \k-0.08) {};
        \node[x_latent, \xCol] at (\i, \j, \k+0.08) {};
      }
    }
  }
\end{tikzpicture}
\caption{The coupled spatio-sequential model as a Forney-style factor graph.
Each spatial node $n$ carries a class label $\ell_n \in \{1, \dots, M\}$
(circles on the lattice, coloured by class for $M = 2$); each class $m$ carries
a hidden chain $k_{s,m} \in \{1, \dots, K\}$ with initial distribution $\Pi_m$
(the two chains at the sides). Pairwise factors (squares) carry the Potts
prior on the labels, Equation~\ref{eq:spatial-prior}, and the chain prior,
Equation~\ref{eq:circulant}. The gated emission factors (triangles) join one
label to one chain state and contribute $\log p(x_{sn} \mid k_{s,m}, \theta_m)$
only when $\ell_n = m$; the observations $x_{sn}$ are the small circles along
each column. Every second observation column is drawn along each spatial axis,
and an illustrative subset of the gated factors, so the structure is legible.}
\label{fig:coupled:sketch}
\end{figure}
